\chapter{Inférence pour la RLS}
\label{sec:InferenceRLS}

\section{Inférence sur les paramètres}
\label{sec:InférenceParamètres}

On a montré comment obtenir des estimations pour les différents paramètres du modèle de régression linéaire simple. On souhaite en général également faire de l'inférence sur les vraies valeurs des paramètres $\beta_0$, $\beta_1$ et $\sigma^2$. Pour ce faire, on utilisera la distribution des estimateurs, que nous dérivons ici.

\subsection{Distribution des estimateurs}
\label{sec:LoisBetas}

\subsubsection{Distribution de $\hat{\beta}_1$}
\label{sec:LoiDeHatBeta1}

%Notez que $\hat{\beta}_1$ est une variable aléatoire puisque c'est une fonction des $Y_i$ qui sont elles-mêmes des variables aléatoires. On veut trouver sa distribution.

On note d'abord que $\hat{\beta}_1$ constitue une combinaison linéaire des $Y_i$. En effet,
$$ \hat{\beta}_1= \frac{\sum_{i=1}^n (x_i-\overline{x})Y_i}{\sum_{i=1}^n (x_i-\overline{x})^2} = \sum_{i=1}^n a_i Y_i
$$
o\`u $a_i= \{(x_i-\overline{x})/\sum_{i=1}^n(x_i-\overline{x})^2\} = (x_i-\overline{x})/S_{xx}$.

Étant donné qu'une combinaison linéaire de variables aléatoires normales suit également une distribution normale, il s'ensuit que $\hat\beta_1$ est distribué selon une loi normale. Il ne reste qu'à trouver l'espérance et la variance de celle-ci.

Notons d'abord que
 \begin{eqnarray*}
 \sum_{i=1}^n a_i & = & \frac{\sum_{i=1}^n(x_i-\overline{x})} {S_{xx}} = 0 \\
\sum_{i=1}^n a_i x_i & = & \frac{\sum_{i=1}^n(x_i-\overline{x})x_i } {S_{xx}} =\frac{\sum_{i=1}^n(x_i-\overline{x})(x_i-\overline{x})}{S_{xx}}=\frac{S_{xx}}{S_{xx}}=1.
\end{eqnarray*}




Pour l'espérance, on peut montrer que $\hat{\beta}_1$ est un estimateur sans biais de $\beta_1$:
\begin{eqnarray*}
 \esp[\hat{\beta}_1] &=& \sum_{i=1}^n a_i \esp[Y_i] \\
 &=& \sum_{i=1}^n a_i (\beta_0 + \beta_1 x_i)\\
 &=& \beta_0 \sum_{i=1}^n a_i + \beta_1 \sum_{i=1}^n a_i x_i\\
 &=& \beta_1
\end{eqnarray*}


Pour la variance de $\hat\beta_1$, on peut calculer
\begin{eqnarray*}
 \var[\hat\beta_1]&=&\var\left[\sum_{i=1}^n a_i Y_i\right]\\
 &=& \sum_{i=1}^n a_i ^2 \var[Y_i]\\
 &=& \sum_{i=1}^n \frac{(x_i-\overline{x})^2}{S_{xx}^2} \var[Y_i]\\
 &=& \frac{1}{S^2_{xx}} \sum_{i=1}^n (x_i-\overline{x})^2 \var[Y_i]\\
 &=& \frac{S_{xx}}{S_{xx}^2} \sigma^2 \\
 &=& \frac{\sigma^2}{S_{xx}}
\end{eqnarray*}

On obtient donc la distribution de $\hat\beta_1$ suivante
\begin{eqnarray}
\label{distbeta1}
 \hat{\beta}_1 \sim N\left( \beta_1, \frac{\sigma^2}{S_{xx}} \right).
 \label{loibeta1}
\end{eqnarray}

\vspace{0.1in}
\begin{remarque}
Seulement certaines hypothèses de la RLS ont été nécessaires pour calculer cette distribution. Lesquelles? Et que nous indique la forme de cette distribution sur l'estimation de la pente dans un modèle de régression linéaire simple?
\end{remarque}


\subsubsection{Distribution de $\hat\beta_0$}
\label{DistributionBeta0}
L'estimateur $\hat\beta_0$ peut également s'exprimer comme une combinaison linéaire des $Y_i$ :
\begin{eqnarray*}
 \hat{\beta}_0 &=& \overline{Y} - \hat{\beta}_1 \overline{x} \\
 &=& \sum_{i=1}^n Y_i /n - \overline{x} \sum_{i=1}^n a_i Y_i \\
 &=& \sum_{i=1}^n (1/n - a_i \overline{x}) Y_i
\end{eqnarray*}
Par le même raisonnement que pour $\hat\beta_1$, on a donc que $\hat\beta_0$ est distribué selon une loi normale. On peut calculer la moyenne et la variance de cette loi en utilisant l'espérance et la variance de $\hat\beta_1$. Pour l'espérance, on a que
\begin{eqnarray*}
 \esp[\hat{\beta}_0] &=& \esp[\overline{Y} - \hat{\beta}_1 \overline{x}] \\
 &=& \frac{1}{n}\sum_{i=1}^n \esp[Y_i] - \overline{x} \esp[\hat{\beta}_1]  \\
 &=& \frac{1}{n}\sum_{i=1}^n(\beta_0+\beta_1 x_i)-\overline{x}\beta_1 \\
 & = & \beta_0+\beta_1 \overline x-\beta_1\overline x \\
 & = & \beta_0
\end{eqnarray*}
La variance s'exprime quant à elle de la façon suivante
\begin{eqnarray*}
 \var[\hat{\beta}_0] &=& \var[\overline{Y} - \hat{\beta}_1 \overline{x}] \\
 &=& \var\left[\frac{1}{n}\sum_{i=1}^n Y_i\right]+ \overline{x}^2 \var[\hat{\beta}_1]-2\overline x\cov[\overline Y,\hat\beta_1]  \\
 &=& \frac{\sigma^2}{n}+\overline{x}^2\frac{\sigma^2}{S_{xx}} \\
 & = & \sigma^2\left(\frac{1}{n}+\frac{\overline{x}^2}{S_{xx}}\right)
\end{eqnarray*}
Pour obtenir l'expression finale pour cette variance, il faut démontrer que $\cov[\overline Y,\hat\beta_1]=~0$, ce qui est laissé en exercice. \\

On obtient ainsi la loi suivante pour $\hat\beta_0$:
\begin{eqnarray}
\label{distbeta0}
 \hat{\beta_0} \sim \N\left\{\beta_0, \left(\frac{1}{n}+\frac{\overline{x}^2}{S_{xx}}\right)\sigma^2\right\}.
\label{loibeta0}
\end{eqnarray}

\begin{exemple}
Calculons la variance de $\hat\beta_0$ et $\hat\beta_1$ dans l'exemple de la section \ref{exemple-RLS-interpretation}. On doit pour ce faire connaître la valeur de la variance des termes d'erreur. Supposons ici que $\sigma^2=1$. On sait déjà que $n=10,\overline x=145$ et $S_{xx}=8250$. Par conséquent on a
$$\var[\hat\beta_0]=\left( \frac{1}{10}+\frac{145^2}{8250} \right)(1)=2,648$$
$$\var[\hat\beta_1]=\frac{1^2}{8250}=1,212\times 10^{-4}.$$
\end{exemple}



\subsection{Intervalles de confiance pour $\beta_0$ et $\beta_1$}
\label{sec:IntervallesDeConfiance}

On peut utiliser les distributions de $\hat{\beta_0}$ et $\hat{\beta_1}$ pour construire des intervalles de confiance pour $\beta_0$ et $\beta_1$ avec la méthode des pivots. Par exemple, pour un intervalle de confiance pour $\beta_1$, on centre et réduit la variable aléatoire pour obtenir
\begin{equation}
\frac{\hat{\beta}_1-\beta_1}{\sqrt{\sigma^2/S_{xx}}}\sim \N(0,1).
\label{disbeta0std}
\end{equation}
%équations \eqref{distbeta0} et \eqref{distbeta1})
\vspace{-0.05 in}
Si $\sigma^2$ était connue, un intervalle de confiance de niveau $100(1-\alpha)\%$ serait donc donné par $$[\hat{\beta_1}\pm z_{1-\alpha/2} \sigma / \sqrt{S_{xx}}].$$ Or, $\sigma^2$ n'est généralement pas connue et on doit par conséquent l'estimer. L'estimateur du maximum de vraisemblance de cette quantité est $\sum_{i=1}^n \hat{\varepsilon}_i^2 /n$ mais ce dernier est biaisé.
%Montrer que c'est le cas.
On utilisera donc plutôt l'estimateur non biaisé défini par $$\hat{\sigma}^2=  \frac{1}{n-2} \sum_{i=1}^n \hat{\varepsilon}_i^2.$$

\vspace{-0.05 in}
Remplacer $\sigma^2$ par $\hat\sigma^2$ dans le terme de gauche de \ref{disbeta0std} a toutefois pour conséquence que cette variable aléatoire ne suit plus une loi normale. On peut néanmoins montrer, en utilisant des notions de statistique mathématique, que
\begin{equation}
 \frac{\frac{\hat{\beta}_1-\beta_1}{\sigma/\sqrt{S_{xx}}}}
 {  \sqrt{\frac{\hat\sigma^2}{\sigma^2} } } \sim \mathcal{T}_{n-2}.
\label{disbeta0stds}
\end{equation}


On peut alors déduire un intervalle de confiance de niveau $100(1-\alpha)\%$ pour $\beta_1$:
\begin{eqnarray}
\left[\hat{\beta}_1 \pm t_{n-2; 1-\alpha/2} \left(\frac{\hat\sigma}{\sqrt{S_{xx}}}\right)\right],
\end{eqnarray}
où $t_{n-2; 1-\alpha/2}$ dénote le $(1-\alpha/2)$-ième quantile d'une loi $\mathcal{T}$ à $(n-2)$ degrés de liberté.\\

 Il est possible de construire de la même manière un intervalle de confiance pour $\beta_0$. On obtient alors l'intervalle de confiance suivant:
 \begin{eqnarray}
\left[\hat{\beta}_0\pm t_{n-2; 1-\alpha/2}\; \hat\sigma\sqrt{{1\over n}+{\overline{x}^2\over S_{xx}}}\right].
\label{ci:beta0}
\end{eqnarray}

\vspace{0.05in}
\begin{remarque}
Pour trouver \ref{disbeta0stds}, on note que
\begin{eqnarray*}
 \frac{\frac{\hat{\beta}_1-\beta_1}{\sigma/\sqrt{S_{xx}}}}{\sqrt{\frac{\sum_{i=1}^n \hat{\varepsilon}_i^2}{\sigma^2 (n-2)}}}=\frac{U}{\sqrt{V/\nu}}
\end{eqnarray*}
où $U=\frac{\hat{\beta}_1-\beta_1}{\sigma/\sqrt{S_{xx}}}$ est distribué selon une loi normale centrée réduite et $V=\frac{\sum_{i=1}^n \hat{\varepsilon}_i^2}{\sigma^2}$ suit une loi du khi-carré à $\nu=(n-2)$ degrés de liberté, puisque c'est une somme du carré de variables aléatoires normales.
Comme il s'agit d'un rapport entre une variable aléatoire normale centrée réduite et la racine carrée d'une variable khi-carré divisée par son degré de liberté  $(n-2)$, on obtient la distribution d'une loi $\mathcal{T}$ de Student à $(n-2)$ degrés de liberté ($\mathcal{T}_{n-2}$):
\begin{eqnarray*}
 \frac{\frac{\hat{\beta}_1-\beta_1}{\sigma/\sqrt{S_{xx}}}}{\sqrt{\frac{\sum_{i=1}^n \hat{\varepsilon}_i^2}{\sigma^2 (n-2)}}} \sim \frac{\N(0,1)}{\sqrt{\chi_{n-2}^2/(n-2)}} \sim \mathcal{T}_{n-2}
 \end{eqnarray*}
\end{remarque}

\vspace{0.1in}
\begin{exemple}
\label{exp:int_beta}
Calculons un intervalle de confiance pour $\beta_0$ et $\beta_1$ avec les données de la section \ref{exemple-RLS-interpretation}. La plupart des quantités nécessaires ont déjà été calculées. Il ne manque que
$$\hat\sigma=\sqrt{\frac{1}{10-2} \sum_{i=1}^{10}\hat\varepsilon_i^2}=\sqrt{\frac{7,224}{8}}=0,950.$$
Si on fixe $\alpha=5\%$, il faut calculer $t_{n-2; 0,975}$. On peut trouver dans une table que ce quantile est  $t_{8;0,975}=2,306$. On obtient alors les intervalles suivants:
$$[\hat{\beta}_0\pm 2,306\times0,950\sqrt{1/10+145^2/8250}]=[\hat{\beta}_0\pm 3,566]=[-6,305;0,827]$$
$$[\hat{\beta}_1\pm 2,306\times 0,950/\sqrt{8250}]=[\hat{\beta}_1\pm 0,024]=[0,459;0,507]$$

\end{exemple}

\newpage
\subsection{Tests d'hypothèses}

On peut également se baser sur les distributions trouvées à la section \ref{sec:LoisBetas}
pour mettre en {\oe}uvre des tests d'hypothèse sur la valeur des coefficients.

L'hypothèse la plus souvent testée concerne la nullité de la pente, à savoir $\mathcal{H}_0: \beta_1 = 0$. Sous $\mathcal{H}_0$, $\frac{\hat{\beta}_1}{\hat\sigma/\sqrt{S_{xx}}} \sim \mathcal{T}_{n-2}$ et $\prob(|t_{n-2}|>|t_{\text{obs}}|)$, où $t_{\text{obs}}$ désigne la statistique du test calculée, nous livre la $p$-valeur du test bilatéral. Si la $p$-valeur est plus petite que le seuil $\alpha$ choisi, on rejette l'hypothèse nulle que $\beta_1=0$. Au contraire, si la $p$-valeur est supérieure au seuil, on ne peut pas rejeter l'hypothèse nulle que $\beta_1=0$.


 Dans le cas où l'hypothèse nulle que l'ordonnée à l'origine $\beta_0=0$ ne peut être rejetée, le modèle peut se simplifier :  $Y_i = \beta_1 x_i + \varepsilon_i$.

\begin{exemple}
Toujours avec les données de la section \ref{exemple-RLS-interpretation}, on trouve
$$ t_{obs} = \frac{\hat{\beta_1}}{\hat\sigma/\sqrt{S_{xx}}} = \frac{0,483}{0,950/\sqrt{8250}} = 46,17 $$
Un logiciel nous permet d'obtenir $5.35 \times 10^{-11}$ comme valeur de $p$ correspondante. On rejette donc l'hypothèse d'une pente nulle. Notez qu'on aurait pu arriver à la même conclusion simplement à l'aide de l'intervalle de confiance construit à l'exemple \ref{exp:int_beta}. En effet, l'intervalle de confiance à 95\% pour $\beta_1$ ne contient pas la valeur de $0$, ce qui permet de rejeter l'hypothèse que $\beta_1 = 0$ au seuil $\alpha = 0.05$.

%Il est à souligner que comme 0 appartient à l'intervalle de confiance pour $\beta_0$, on ne peut pas rejeter l'hypothèse nulle que l'ordonnée à l'origine soit égale à zéro. Le modèle peut donc être simplifié.
\end{exemple}


\section{Prévision}
\label{sec:PrévisionRLS}

On a déjà défini $\hat{Y}_i = \hat{\beta}_0 + \hat{\beta}_1 x_i$, la prévision du modèle en $x_i$. On souhaite maintenant obtenir un intervalle de confiance pour cette valeur.

On doit alors distinguer deux cas:
\begin{enumerate}
\item[(i)] Estimer $\esp[Y|x=x_0]$, la valeur moyenne de $Y$ prédite par le modèle pour une valeur $x=x_0$
de la variable exog\`ene.
\item[(ii)] Estimer la valeur d'une observation/r\'ealisation de $Y$ lorsque
la variable exog\`ene prend la valeur $x_0$. Cette valeur sera dénotée par ${Y}_0^*$.
\end{enumerate}

Dans le premier cas, on souhaite estimer la valeur de la droite en $x_0$, alors que dans le deuxième cas on s'intéresse à une réalisation de la variable dépendante, qui ne sera pas nécessairement sur la droite. L'intervalle de confiance sera plus grand dans le deuxième cas, pour tenir compte de la variabilité additionnelle causée par la fluctuation aléatoire autour de la droite.
%Les détails sont présentés plus bas.

\subsection{Prévision de la valeur moyenne $\esp[Y|x=x_0]$}

Soit $\esp[Y|x=x_0]$, la valeur moyenne de $Y$ prédite par le modèle lorsque $x=x_0$. Cette quantité est estimée par $\hat{Y}_0 = \hat{\beta}_0 + \hat{\beta}_1 x_0$.
On utilise le fait que $Y_i \sim \N(\beta_0+\beta_1x_i,\sigma^2)$ pour trouver la distribution de $\hat{Y}_0$.
On exprime d'abord $\hat{Y}_0$ sous la forme d'une combinaison linéaire des $Y_i$ de la façon suivante:
\begin{eqnarray*}
 \hat{Y}_0 &=& \hat{\beta}_0 + \hat{\beta}_1 x_0 \\
 &=&\sum_{i=1}^n (1/n - a_i \overline{x}) Y_i + \left(\sum_{i=1}^n a_i Y_i\right) x_0 \\
 &=& \sum_{i=1}^n (1/n - a_i \overline{x} + a_i x_0) Y_i\\
 &=& \sum_{i=1}^n c_i Y_i
\end{eqnarray*}
où
 \[c_i = 1/n - a_i \overline{x} + a_i x_0.\]

On peut montrer que $\sum_{i=1}^n c_i = 1$, que $\sum_{i=1}^n c_i x_i = x_0$ et que $\sum_{i=1}^n c_i^2 = \frac{1}{n} + \left\{\frac{(\overline{x}-x_0)^2}{S_{xx}}\right\}$.

On peut ensuite en déduire que
\begin{eqnarray}
 \esp[\hat{Y}_0]= \sum_{i=1}^n c_i \esp[Y_i]= \sum_{i=1}^n c_i (\beta_0 + \beta_1 x_i) = \beta_0 + \beta_1 x_0 \\
 \label{eq:varyo}
 \var[\hat{Y}_0]= \sum_{i=1}^n c_i^2 \var[Y_i] = \left\{ \frac{1}{n} + \frac{(\overline{x}-x_0)^2}{S_{xx}}\right\} \sigma^2
\end{eqnarray}

Puisque $\hat{Y}_0$ est une combinaison linéaire des $Y_i$, qui suivent une loi normale, $\hat{Y}_0$ suit lui-même une loi normale. On a donc

\begin{equation}
\hat{Y}_0 \sim \N\left(\beta_0 + \beta_1 x_0, \left\{ 1/n + \frac{(\overline{x}-x_0)^2}{S_{xx}}\right\} \sigma^2 \right).
\end{equation}

À partir de ce résultat, il est possible de calculer un intervalle de confiance autour de la valeur moyenne prédite en remplaçant $\sigma^2$ par son estimateur $\hat\sigma^2$. En utilisant le même raisonnement qu'à la section \ref{sec:IntervallesDeConfiance}, on obtient l'intervalle de confiance de niveau $1-\alpha$ suivant pour $\esp[Y|x=x_{0}]$:
\begin{eqnarray}
\left[(\hat{\beta}_0 + \hat{\beta}_1 x_0)\; \pm\; t_{n-2; 1-\alpha/2} \;\hat\sigma \; \sqrt{\frac{1}{n}+\frac{(\overline x-x_0)^2}{S_{xx}}}\right]
\label{eq:int_conf_val_moy_prev}
\end{eqnarray}

 \subsection{Prévision d'une observation $Y_0^*$ lorsque $x=x_0$}

Il est important de comprendre que l'intervalle de confiance défini à l'équation \eqref{eq:int_conf_val_moy_prev} concerne la moyenne de $Y$ au point $x=x_0$. Pour construire ce qu'on appelle un intervalle de prévision, il faut considérer à la fois la variance liée au modèle (à l'estimation des paramètres) et la variance liée à la variabilité naturelle du phénomène, c'est-à-dire la variance du terme d'erreur ($\varepsilon$). Or, l'équation \eqref{eq:varyo} ne tient pas compte de la variance du terme d'erreur, qui est égale à $\sigma^2$.

Désignons la prévision de la moyenne de $Y$ en  $x_0$ par $\hat{Y}_0 = \hat{\beta}_0 + \hat{\beta}_1 x_0$ et la prévision d'une observation individuelle par $\hat{Y}_0^* = \hat{\beta}_0 + \hat{\beta}_1 x_0 + \varepsilon$. L'espérance de $\hat{Y}_0$ et celle de $\hat{Y}_0^*$ sont égales parce que l'espérance des termes d'erreur est nulle. Ce n'est pas le cas des variances.

On trouve en effet que
\begin{eqnarray}
 \var[\hat{Y}_0^*] &=& \var[\hat{Y}_0] + \var [\varepsilon] \nonumber\\
 &=&\left\{ \frac{1}{n} + \frac{(\overline{x}-x_0)^2}{S_{xx}}\right\} \sigma^2 + \sigma^2 \nonumber\\
 &=& \left\{1+\frac{1}{n} + \frac{(\overline{x}-x_0)^2}{S_{xx}} \right\} \sigma^2.
 \label{eq:varyo*}
\end{eqnarray}

Le méthode de calcul de l'intervalle de prévision de niveau $1-\alpha$ pour ${Y}_0^*$ est similaire à celle de l'intervalle de confiance pour  $\esp[Y|x=x_0]$. On obtient donc
\begin{eqnarray}
\left[(\hat{\beta}_0 + \hat{\beta}_1 x_0)\; \pm\; t_{n-2; 1-\alpha/2} \; \hat\sigma \; \sqrt{1+\frac{1}{n}+\frac{(\overline x-x_0)^2}{S_{xx}}} \right].
\label{eq:int_prev}
\end{eqnarray}


En analysant les formules \eqref{eq:varyo} et \eqref{eq:varyo*}, nous constatons que les variances diminuent lorsque $n$ augmente ou que $S_{xx}$ augmente. Cela signifie que plus la taille de l'échantillon est grand et/ou plus il y a de la variabilité dans les valeurs de la variable exogène, plus les estimations seront précises. Au contraire, plus le point $x_0$ pour lequel une prévision est requise est éloigné du centre $(\overline x)$ des variables explicatives, plus la prévision sera imprécise.
Il est important de prendre ces remarques en considération lors de la planification des expériences.


\begin{remarque}
L'intervalle de prévision pour $Y_0^*$ sera toujours plus large que l'intervalle de confiance pour $\esp[Y|x=x_0]$. Pour cette raison, il peut arriver dans la pratique qu'un client préfère présenter l'intervalle de confiance plutôt que l'intervalle de prévision. Il est important de bien comprendre et de bien expliquer la distinction entre les deux intervalles pour éviter une mauvaise interprétation des résultats.
\end{remarque}


\begin{exemple}
On désire calculer un intervalle de confiance à 95 $\%$ pour le rendement moyen du produit ($\esp[Y|x=x_0]$) lorsque la température de fonctionnement est de 150 $^{\circ}\mathrm{C}$ et pour le rendement d'un produit spécifique (${Y}_0^*$) lorsque la température est de 150 $^{\circ}\mathrm{C}$.

Dans les deux cas, la prévision est $\hat\beta_0+ \hat\beta_1\times150 =-2,739+0,483\times 150=69,711 \%$.

La plupart des quantités nécessaires au calcul des deux intervalles ont été déjà calculées auparavant. On tire de la section \ref{exemple-RLS-interpretation}  que $\overline x=145, S_{xx}=8250$ et de l'exemple \ref{exp:int_beta} que $\hat\sigma=0,950$ et que $t_{8; 0,975}=2,306$. Il nous reste seulement à calculer $(x_0-\overline x)^2=(150-145)^2=25$.

En utilisant l'équation \eqref{eq:int_conf_val_moy_prev}, on obtient l'intervalle de confiance du rendement moyen :
$$[69,711\pm 2,306\times 0,950\times\sqrt{1/10+25/8250}]=[69,711 \pm 0,703]=[69,008;70,414].$$

L'intervalle de prévision du rendement d'un produit quelconque, appelé également intervalle de prévision, se calcule par l'équation~\eqref{eq:int_prev}: \\
 $[69,711\pm 2,306\times 0,950\times\sqrt{1+1/10+25/8250}]=[69,711\pm 2,301]=[67,410;72,012]$.

Comme on l'a déjà signalé l'intervalle de prévision est plus large que l'intervalle de confiance pour $\esp[Y|x=150]$.
\end{exemple}


\section{Analyse de la variance}
\label{sec:AnalyseDeLaVariance}

L'analyse de la variance est vue en profondeur dans un cours qui lui est dédié, mais c'est un outil qui peut également être utile dans une analyse de régression. On verra dans cette section comment la variabilité dans les observations de la variable endogène $Y_1,\ldots,Y_n$ peut se décomposer sous le modèle de régression linéaire. Cette décomposition constitue la base des tests d'ajustement du modèle.

\subsection{Décomposition de la variabilité}

On note d'abord que les valeurs de la variable réponse ne seront pas toutes identiques dans le jeu de données. Une partie de la différence entre les $Y_i$ peut être expliquée par la différence entre les $x_i$, puisque selon le modèle de régression la variable réponse est expliquée en partie par la variable explicative. Il restera toutefois une partie de la variabilité des $Y_i$ qui n'est pas expliquée par les différences dans les valeurs de $x_i$. Mathématiquement, on peut écrire
\begin{align*}
Y_i - \overline{Y} &= (\hat{Y}_i - \overline{Y})  +  (Y_i - \hat{Y}_i)
\end{align*}
La base de l'analyse de variance en régression, c'est que la même relation est vraie au niveau de la somme des carrés totaux:
\begin{eqnarray*}
  \sum_{i=1}^n (Y_i-\overline{Y})^2 & = & \sum_{i=1}^n (Y_i - \hat{Y}_i + \hat{Y}_i - \overline{Y})^2\\
 &=&  \sum_{i=1}^n (Y_i - \hat{Y}_i)^2 + \sum_{i=1}^n(\hat{Y}_i - \overline{Y})^2 + 2\sum_{i=1}^n (Y_i-\hat{Y}_i)(\hat{Y}_i-\overline{Y})\\
 &=&  \sum_{i=1}^n (Y_i - \hat{Y}_i)^2 + \sum_{i=1}^n(\hat{Y}_i - \overline{Y})^2.
\end{eqnarray*}

\vspace{-0.05 in}
On écrira plus simplement
\begin{equation}\label{eq:decompvar}
SC_{tot}=SC_{reg}+SC_{res}.
\end{equation}
o\`u
\vspace{-0.05 in}
\begin{align*}
SC_{tot}& =\sum_{i=1}^n(Y_i-\overline{Y})^2 \; , \mbox{ la somme des carrés totale}, \\
SC_{reg} & =\sum_{i=1}^n(\hat{Y}_i-\overline{Y})^2 \; \mbox{ la somme des carrés de la régression, et}\\
SC_{res}&=\sum_{i=1}^n(Y_i-\hat{Y}_i)^2 \; \mbox{ la somme des carrés résiduelle.}
\end{align*}
La somme des carr\'es totaux $SC_{tot}$ quantifie la variabilit\'e dans les $Y_i$ autour de leur valeur moyenne.  La somme des carrés de la régression $SC_{reg}$ quantifie la variabilit\'e dans les prévisions $\hat{Y}_i$. Étant donné que les $\hat{Y}_i$ ne varient qu'avec les $x_i$, cette somme de carrés est la partie de la variabilit\'e dans les $Y_i$ expliqu\'ee par le fait que toutes les observations n'ont pas la m\^eme valeur pour $x_i$. La somme  des carrés résiduelle $SC_{res}$ mesure la variabilit\'e dans les $(Y_i-\hat{Y}_i)$. Cette variabilit\'e est induite par le fait que la valeur de $Y_i$ n'est pas expliqu\'ee parfaitement par le mod\`ele de
r\'egression.

%Cette variabilit\'e se d\'ecompose en deux parties : la somme des carrés de la régression ($SC_{reg}$) et la somme des carrés résiduelle ($SC_{res})$.

\begin{exemple}
\label{exp:ss}
Toujours dans l'exemple du rendement d'un produit en fonction de la température de fonctionnement du procédé, on peut calculer que
{\small
\begin{align*}
SC_{tot} &= \sum_{i=1}^n(Y_i-\overline{Y})^2 = 1931,2 \\
SC_{res} &= \sum_{i=1}^n(Y_i-\hat{Y}_i)^2 = (n-2)*\hat\sigma^2 = 7,224 \\
SC_{reg} &= \sum_{i=1}^n(\hat{Y}_i-\overline{Y})^2 = 1924,876
\end{align*}
}
On vérifie alors que $SC_{tot} = SC_{reg} + SC_{res}$. On note que la température \`a elle seule ne peut pas expliquer en totalit\'e pourquoi toutes les expériences ne possèdent pas le m\^eme rendement.
\end{exemple}

\subsection{Test global de la régression}

La somme des carrés de la régression ($SC_{reg}$) constitue donc la portion de la variance expliquée par le modèle, alors que la somme des carrés résiduelle ($SC_{res}$) représente la partie que le modèle n'explique pas. Intuitivement, on comprend que si la proportion de la variance expliquée par le modèle ($SC_{reg}/SC_{tot}$) est élevée, le modèle est bon.

L'analyse de la variance utilise cette idée pour effectuer un test global de la régression. On désire tester les hypothèses suivantes
\begin{align}
H_0&:Y_i=\beta_0+\varepsilon_i \nonumber \\
H_1&:Y_i=\beta_0+\beta_1x+\varepsilon_i  \label{eq:HypotheseAnova}
\end{align}
L'idée consiste à comparer $SC_{reg}$ à $SC_{res}$ pour évaluer si la variable exogène explique une partie significative de $SC_{tot}$. Afin de construire un test, il faut en premier lieu trouver la distribution de la statistique sous $H_0$. Il se trouve qu'il est plus simple de trouver cette distribution si on standardise les sommes des carrés pour construire la statistique de test. On obtient alors $CM_{res}$, le carré moyen pour la régression, et $CM_{reg}$, le carré moyen pour les résidus. La statistique utilisée pour tester les hypothèses \ref{eq:HypotheseAnova} sera le rapport des carrés moyens $F = CM_{reg}/CM_{res}$. Le \autoref{tab:tabanova} présente les principaux éléments de l'analyse de la variance et montre comment calculer cette statistique de test.

\begin{table}[htb!]
\caption{Tableau d'analyse de la variance dans le cas de la régression linéaire simple présentant la source, le nombre de degrés de liberté, la somme des carrés, le carré moyen et la statistique $F$ de Fisher.}
\centering
\begin{footnotesize}
\renewcommand{\arraystretch}{1.6}
\begin{tabular}{|c|c|c|c|c|}\hline
Source & Degrés de liberté & Somme des carrés & Carrés moyens & $F$ \\ \hline
Modèle & $ddl_1=p=1$ & $SC_{reg}=\sum_{i=1}^n (\hat Y_i - \overline Y)^2$ & $CM_{reg}=SC_{reg}/ddl_1$  & $CM_{reg}/CM_{res}$ \\ \hline
Erreur & $ddl_2=n-2$ & $SC_{res}=\sum_{i=1}^n (Y_i - \hat Y_i)^2$ & $CM_{res}=SC_{res}/ddl_2$ & \vspace*{-3 mm}\\
Résiduelle & & & & \\ \hline
Totale & $n-1=n-2+p$ & $SC_{tot}=\sum_{i=1}^n (Y_i - \overline Y)^2$ &  & \\ \hline
\end{tabular}
\end{footnotesize}
\label{tab:tabanova}
\end{table}

Une grande valeur de la statistique $F$ indique que la régression explique une bonne partie de la variabilité, et suggère de rejeter l'hypothèse nulle. Pour calculer la zone de rejet, il faut utiliser la distribution de la statistique de test sous $H_0$. On trouvera sous peu que $F = CM_{reg} / CM_{res} \sim F_{1,n-2}$.

On rejette donc $H_0$ au seuil $\alpha$ si $F\geq \mathcal{F}_{1,n-2;1-\alpha}$. De plus, une faible valeur de $p$ pour ce test (par exemple si $p<0,05$) indiquera que l'effet de la variable exogène est significatif.
%La $p$-valeur du test de Fisher-Snedecor indique si la variable exogène possède un effet significatif sur la valeur de la variable endogène. Une $p$-valeur faible .

\textbf{Distribution de la statistique sous $H_0$}

On peut trouver que sous $H_0$:
\begin{eqnarray*}
F =  \frac{CM_{reg}}{CM_{res}}&=& \frac{SC_{reg}/1}{\hat\sigma^2}\\
 &=& \frac{S_{xy}^2/S_{xx}}{\hat\sigma^2}\\
 &=& \frac{\hat{\beta}_1^2 S_{xx}}{\hat\sigma^2}\\
 &=& \left(\frac{\hat{\beta}_1}{\hat\sigma/\sqrt{S_{xx}}}\right)^2 \sim (\mathcal{T}_{n-2})^2 \sim \mathcal{F}_{1,n-2}
\end{eqnarray*}
Dans la dernière égalité, nous avons utilisé le fait que si la variable aléatoire $T$ est distribuée selon la loi $\mathcal{T}$ de Student à $k$ degrés de liberté (notée $\mathcal{T}_k$) alors la variable aléatoire $T^2$ suit la loi de Fisher-Snedecor à $1$ et $k$ degrés de liberté (notée $\mathcal{F}_{1,k}$).

%On va donc rejeter $H_0$ au seuil $\alpha$ si $F\geq \mathcal{F}_{1,n-2}(1-\alpha)$. La $p$-valeur du test de Fisher-Snedecor indique si la variable exogène possède un effet significatif sur la valeur de la variable endogène. Une $p$-valeur faible (par exemple si la $p$-valeur $<0,05$) signifie que l'effet de la variable exogène est significatif.

\vspace{0.05in}
\begin{remarque}
Dans le tableau \ref{tab:tabanova} d'analyse de la variance, chaque somme des carrés possède un nombre de degrés de liberté qui lui est propre (indiqué dans la deuxième colonne du tableau), à savoir :
\begin{itemize}
\item[$\bullet$] $SC_{reg}$ possède un seul degré de liberté ($ddl_1=p=1$) car dans le modèle de régression linéaire simple il n'existe qu'une seule variable explicative (exogène).
\item[$\bullet$] $SC_{res}$ est associé à $(n-2)$ degrés de libertés car on estime deux paramètres dans le modèle de régression ($\beta_0$ et $\beta_1$), donc on a $n$ observations moins 2 paramètres.
\item[$\bullet$] $SC_{tot}$ possède $(n-1)$ degrés de liberté car il y a $n$ observations moins une moyenne.
\end{itemize}
On note que le nombre de degrés de liberté de $SC_{tot}$ est égal à la somme du nombre de degrés de liberté de $SC_{reg}$ et $SC_{res}$.
\end{remarque}


\begin{remarque}
Il est à souligner que le carré moyen de l'erreur résiduelle ($CM_{res}$) correspond à l'estimation de la variance des termes d'erreur. En effet
 \[CM_{res}=\frac{SC_{res}}{ddl_2}= \frac{\sum_{i=1}^n (Y_i-\hat{Y}_i)^2}{n-2}=\frac{\sum_{i=1}^n \hat\varepsilon_i^2}{n-2}=\hat\sigma^2.
\]
\end{remarque}

\vspace{0.05 in}
\begin{remarque}
Dans le cadre de la régression linéaire simple, le test de Fisher-Snedecor est exactement égal au test de Student de la nullité de la pente $H_0: \beta_1=0$. On verra que cette remarque ne se généralise pas au cas de la régression linéaire multiple.
\end{remarque}

\newpage
\begin{exemple}
On peut calculer le tableau d'analyse de la variance à partir de nos calculs à l'exemple \ref{exp:ss}. On obtient


\begin{table}[htb!]
\centering
\begin{small}
%\renewcommand{\arraystretch}{1.6}
\begin{tabular}{|r|r|r|r|r|}\hline
Source & Degrés de liberté & Somme des carrés & Carrés moyens & $F$ \\ \hline
Modèle & $ddl_1=1$ & 1924,876 & 1924,876   & 2131,574 \\ \hline
Erreur résiduelle & $ddl_2=8$ &  7,224 & 0,903 & \\ \hline
Totale & $n-1=9$ & 1931,1&  & \\ \hline
\end{tabular}
\end{small}
\label{tab:tabanova_exp}
\end{table}

La $p$-valeur du test de Fisher-Snedecor est donné par $\prob(\mathcal{F}_{1,8}>2131,574)=5,352\times 10^{-11}$. On rejette donc l'hypothèse nulle avec une grande confiance et on en conclut que le modèle de régression linéaire est adapté.

Notez qu'on aurait pu utiliser l'estimateur de la variance des termes d'erreur $\hat\sigma^2=0,903$ de l'exemple \ref{exp:int_beta} pour calculer la somme des carrés de l'erreur résiduelle:
$$SC_{res}=CM_{res}\times ddl_2=0,903\times 8=7,224.$$ On peut également calculer plus simplement $SC_{reg}=SC_{tot}-SC_{res}=1932,1-7,224=1924,876$ au lieu de faire le calcul au long.
%Comme $ddl_1=1$, on a dans ce cas que $CM_{reg}=SC_{reg}=1924,876$. Toutes les quantités nécessaires  ont donc été calculées et on obtient le tableau d'analyse de la variance suivant :

\end{exemple}


\subsection{Coefficient de détermination}
\label{sec:R2}

On peut évaluer la qualité du modèle à l'aide du test de Fisher ou du test de nullité de $\beta_1$. Si la conclusion du test est que $\beta_1 = 0$, le modèle est inadapté, c'est-à-dire que la variable explicative (exogène) ne possède pas d'effet significatif linéaire sur la variable réponse (endogène).

L'analyse de la variance suggère également une façon de quantifier la qualité du modèle, soit le coefficient de détermination:

\begin{equation}
R^2 =\frac{SC_{reg}} {SC_{tot}} \in[0,1]
\end{equation}

Le $R^2$ s'interprète comme la proportion de la variance des $Y_i$ expliquée par le modèle. Une valeur élevée indique que le modèle est bon. Le seuil à partir duquel on considère que $R^2$ est élevé est très variable et assez subjectif. Il dépend notamment des objectifs de la régression et du domaine d'application.

Il se trouve que le coefficient de détermination mesure en fait le carré de la corrélation entre les valeurs observées et les valeurs prédites
 \[R^2 = \frac{SC_{reg}} {SC_{tot}} = \corr^2 (Y, \hat{Y})=\left(\frac{\sum_{i=1}^n (Y_i - \overline{Y})(\hat{Y}_i - \overline{Y})}{\sqrt{\sum_{i=1}^n (Y_i - \overline{Y})^2\sum_{i=1}^n(\hat{Y}_i - \overline{Y})^2}}\right)^2
\]

Quand $R^2=0$, toute la variabilit\'e est due \`a l'erreur al\'eatoire
et le mod\`ele n'explique absolument rien de la valeur de $Y_i$. Quand $R^2=1$, tous les points sont alignés sur la droite de r\'egression, c'est-à-dire que l'ajustement du mod\`ele est parfait et que la valeur de $Y_i$ est une fonction exacte de $x_i$.

\begin{exemple}
Dans notre exemple,
$$R^2 = \frac{1924,876}{1931,1} = 0,9967$$
ce qui indique que la température prédit presque parfaitement le rendement de l'expérience. Des valeurs aussi grandes sont plutôt rares en pratique.
\end{exemple}
